%%%
%%% STAT 5014 - Your first LaTeX document
%%% Compile with:  pdflatex example.tex     (or: latexmk -pdf example.tex)
%%% Everything before \begin{document} is the "preamble" and sets up your document
%%%
\documentclass[11pt]{article}

\usepackage{amsmath}                 % align, cases, \text, \operatorname, ...
\usepackage{amssymb}                 % \mathbb, \mathcal, extra symbols
\usepackage{amsthm}                  % theorem and proof environments
\usepackage[margin=1in]{geometry}    % sensible margins

% Operators: upright and correctly spaced, unlike typing "Var(X)"
\DeclareMathOperator{\Var}{Var}
\DeclareMathOperator{\Cov}{Cov}
\DeclareMathOperator*{\argmin}{arg\,min}   % the * puts limits underneath

% Macros: teach LaTeX your notation
\newcommand{\E}{\mathbb{E}}
\newcommand{\R}{\mathbb{R}}
\newcommand{\iid}{\overset{\text{iid}}{\sim}}
\newcommand{\norm}[1]{\left\lVert #1 \right\rVert}   % takes one argument

% Theorem environments
\newtheorem{theorem}{Theorem}
\theoremstyle{definition}
\newtheorem{definition}{Definition}

\title{A First \LaTeX{} Document}
\author{Your Name}
\date{\today}

\begin{document}
\maketitle

\section{Text and inline math}

This is ordinary text.
A blank line starts a new paragraph; a single line break does not.
Emphasis is \emph{italic} or \textbf{bold}, and code is \texttt{typewriter}.
Note that each sentence starts on a new line in the source: LaTeX doesn't care,
but Git does, and your diffs will thank you.

Quotation marks need ``two backticks and two apostrophes,'' not "straight quotes."
Special characters must be escaped with a backslash: 50\% of the time you need
\%, \$, \&, \_, or \#.
(A bare \% silently comments out the rest of the line, which is a classic first bug.)

Inline math sits between dollar signs: the sample mean is
$\bar{x} = \frac{1}{n}\sum_{i=1}^{n} x_i$, and we often assume $x \in \R$.

\section{Displayed equations}

Use \verb|\[ ... \]| for an unnumbered displayed equation, or the
\texttt{equation} environment for a numbered one:
\begin{equation}
  f(x) = \frac{1}{\sqrt{2\pi}\,\sigma}
         \exp\!\left( -\frac{(x-\mu)^2}{2\sigma^2} \right).
\end{equation}

Multi-line derivations line up with \texttt{align}; the \texttt{\&} marks the
alignment point and \verb|\\| ends a line.
Note \verb|\Var| and \verb|\E| below (those are our own macros):
\begin{align}
  \Var(X) &= \E\!\left[(X-\mu)^2\right] \\
          &= \E[X^2] - \mu^2 .
\end{align}

The ordinary least squares estimator solves a minimization problem:
\[
  \hat{\beta} = \argmin_{\beta \in \R^p} \norm{y - X\beta}^2
              = (X^\top X)^{-1} X^\top y .
\]

\section{Statistical notation}

\begin{itemize}
  \item Distributed as: $X \sim \mathcal{N}(\mu, \sigma^2)$.
  \item Independent and identically distributed: $X_1, \ldots, X_n \iid F$.
  \item Convergence: $\bar{X}_n \xrightarrow{p} \mu$ and
        $\sqrt{n}(\bar{X}_n - \mu) \xrightarrow{d} \mathcal{N}(0, \sigma^2)$.
  \item Independence: $X \perp\!\!\!\perp Y$.
  \item Conditional densities use \verb|\mid|, not a bare bar: $f(y \mid x)$.
  \item Binomial coefficients: $\binom{n}{k}$.
\end{itemize}

\section{Matrices and piecewise functions}

\texttt{pmatrix} gives parentheses (\texttt{bmatrix} gives square brackets):
\[
  \Sigma = \begin{pmatrix}
             \sigma_1^2            & \rho\,\sigma_1\sigma_2 \\
             \rho\,\sigma_1\sigma_2 & \sigma_2^2
           \end{pmatrix}.
\]

\texttt{cases} handles piecewise definitions:
\[
  \mathbf{1}\{x > 0\} =
  \begin{cases}
    1 & \text{if } x > 0, \\
    0 & \text{otherwise.}
  \end{cases}
\]

\section{Theorems and proofs}

\begin{definition}
  An estimator $\hat{\theta}$ is \emph{unbiased} for $\theta$ if
  $\E[\hat{\theta}] = \theta$ for all $\theta$.
\end{definition}

\begin{theorem}[Central Limit Theorem]
  Let $X_1, \ldots, X_n \iid F$ with mean $\mu$ and finite variance $\sigma^2$.
  Then $\sqrt{n}(\bar{X}_n - \mu) \xrightarrow{d} \mathcal{N}(0, \sigma^2)$.
\end{theorem}

\begin{proof}
  Beyond the scope of this example... which is a phrase you should get comfortable
  writing.
\end{proof}

\section{Lists}

Bulleted with \texttt{itemize}:
\begin{itemize}
  \item first point
  \item second point
\end{itemize}

Numbered with \texttt{enumerate}:
\begin{enumerate}
  \item Show that the estimator is unbiased.
  \item Derive its variance.
    \begin{enumerate}
      \item Start from the definition.
    \end{enumerate}
\end{enumerate}

\end{document}
